\cotefile{29}
\scanpage{36}
as developped in Exp VI, and under a mild restriction \inserted{in the nature} of a \unsure{connectedness} assumption, we can say that the category in question is \deleted{\unsure{classifind|classified}} completely described by a certain pro-finite group, the so-called fundamental group of $(Z,G)$, which in fact is an extension of $G$ by the usual fundamental group of $Z$. On the other hand, for these latter considerations, all normality assumptions are completely superflu\deleted{\unsure{s|u}}ous. Their main use was to insure that the functor $Z'\quad Z'/G$ from étale coverings of $(Z,G)$ to coverings of $S$ is fully faithful; this latter fact, on the other hand, reduces at once to the case when $Z$ is a principal covering of $S$, with group $G$ (\unsure{because} this condition will be satisfied when restricting to a non empty open set $U$ of $S$, and on the pther hand the restriction functors \deleted{\unsure{fr}} to $U$ resp. $Z/U$ are fully faithful, because of normality); but in this case descent theory tells us that we get in fact an \uline{equivalence} of the category of \uline{étale} coverings of $S$, and the catégory of étale coverings of $(Z,G)$ !

In most applications however, instead of having just one model for ramification (one $Z_a^n$ say), we have an infinity of such (all $Z_a^n$ say, \unsure{fixed} $a$ and variable $n$), and the condition on coverings \inserted{$S'$} of $S$ we are \unsure{interested} in is that the covering should (at least locally) have ramification no \unsure{wosse|worse} than one of the $Z_i$ . This situation is described in example twp, On the other hand, sometimes the ramification models are not defined globally (for instance in the Kummer case, when we give \deleted{ideals which \unsure{ar|as}} divisors which are not principal), only locally, in a way that among the set of \unsure{modals|models} \unsure{given} over various open sets (or more generally, objects \unsure{of} the étale site) there are certain compatibility conditions, roughly that the restriction of any model over $U$ to $U\;V$ has ramification no \unsure{worse} than the restriction of some model given over $V$. Such an example is
% archive sheet number: 35
\scanpage{37}
dealt with in example 3, which has the disadvantage of being somewhat involved and artificial in its technical details, but has the merit of allowing explicit construction of the functor ``geometric realization''. Otherwise, the set up outlined before the examples seems more natural, and giving a closer grasp of the \deleted{\unsure{main}} really essential features of what ``ramification data'' should really mean.(By the way, one shoul rather say ``category of ramification data'').

You are of course right that in the definition of a Galois covering when I \unsure{wssume|assume} that the \unsure{soverings|coverings} of the various $\underline{s}(M)$ are Galois étale coverings with Galois group $G$, that beforhand an action of $\unsure{\mathbb{E}|E|G}$ on $X$ was given, which means that $G$ acts on each of the coverings $X(M)$ of $\underline{s}(M)$ on a way compatible with the base changes $\underline{s}(M')\quad\underline{s}(M)$ stemming from arrows \deleted{\unsure{along the}} $M'\quad M$. In the various examples 1 to 3, this means that the action of $G$ commutes to \inserted{étale} localization on $S$ (i.e. base changes $\unsure{S''|S'}\quad S'\deleted{\unsure{S}}$ in the étale site of $S$) \uline{and} to the action of the groups of automorphisms given on the various $Z$'s .

Now let's come back to the case $S$ strictly local; then (for \unsure{coverings|coveings} defined over t he whole of $S$) the general situation of example 3 reduces to the one of example 2, and replacing the $Z_i$'s by suitable connected components, we may assume them connected if we want. For a given $\unsure{Z_i|Z_i'}$, the $\deleted{\unsure{f_i}}$ étale coverings of $(Z_i,G_i)\deleted{\unsure{x|,}}$ are trivial as coverings of $Z_i$ alone (i.e. forgetting the action of $G_i$) be\deleted{\unsure{r|x}}cause $Z_i$ is strictly local too. Therefore the catégory of étale coverings of$(Z_i,\inserted{G_i)}$ is just equivalent to the category of finite sets $E$ on which $G_i$ opérates; to such an $E$ is associated the consatnt covering $E_{Z_i}$ of $Z_i$, byt $G_i$ operating on it via it's operation on $E$; the geometric realization over $S$ of this is nothing else but $Z_i\overset{G}{x}E$ , the \unsure{vovering|covering|wovering} associated to the $G$-covering
% archive sheet number: 36
\scanpage{38}
$Z_i$ of $S$ and the action of $G_i$ on $E$.\inserted{When $S$ is again arbitrary,}This shows in a \unsure{rahher|rather} concrete way what it means, \inserted{locally(for étale topology)} for a covering $S'$ of $S$, to be associated to some $M$: it means that locally it is describable by a representation of an inertia\deleted{\unsure{l}} group on a finite set, in the obvious way; and the extra structure on one $S'$ involved by deducing it from some $S$, can be described by saying that we give such local representations of $S'$, in such a way that they match together i.e. satisfy some rather evident compatibility condition (This description makes a sense whenever the functor ``\unsure{geometric} realization'' is faithful - not necessarily fully faithful). The analogous description holds for Galois coverings with Galois group $\unsure{\mathbb{G}|G|H}$, \unsure{namely} in the strictly local case, for a fixed $Z_i$ , the\deleted{y are \unsure{classified} by \unsure{systems} (} category of these is equivalent to the category of finite sets $E$ on which $G_i$ operates on the left, $H$ on the right, $G_i$ and $H$ commuting, in such a way that $E$ is a principal homogeneous space (or, as we say, a torsor) under $H$; up to isomorphism, they are classified by $\operatorname{Hom}(G_i,H)/H$ , group homomorphisms mod. interior automorphism. Using \deleted{\unsure{th}} this description, one easily proves the criterion I indicated for the ``geometric\deleted{\unsure{r|-}}\unsure{reallization|realization} -functor'' to be fully faithful on the category of \unsure{\c{\u{G}}|G}alois coverings of type $\underline{R}$ ; indeed, as usual the question is a \unsure{pocal|local} one for the étale topology, which allows one to \deleted{\unsure{restr}} reduce to\unsure{'} the stricly local case.

Your interpretation of the Kummer case in the final \unsure{sormulation|formulation} of example 3 is indeed the one I had in \unsure{mihd|mind}. Also, when I wrote $n'\quad n$, I meant of course the order relation of divisibility (it may be convevient to introduce this order relation explicitly, for simplicity of notations).
% archive sheet number: 37
\scanpage{39}

I realize that all the indications I have given you so far are extremely sketshy, and as a consequence that I am charging you with a considerable amount of work to put some sense and order into all that. Thus it is I, not you, who should apologize for causing a lot of trouble! I look forward with great \unsure{pladsure|pleadsure|pleasure} meeting you in Bures. As I am having some russian andchinese lessons on fridays, I will probably drop by on June 2.

With best regards\\

\handwritten{\figure{\begin{tikzpicture}[x=1mm,y=1mm]
\draw[line width=.35mm,gray] (0,2.6) .. controls (.1,1.8) and (.3,.8) .. (.65,.15) .. controls (.95,-.05) and (1.1,.1) .. (.8,.75) -- (.3,2.6);
\draw[line width=.35mm,gray] (.6,2.5) .. controls (.9,1.6) and (1.1,.7) .. (1.4,.2) .. controls (1.65,-.05) and (1.7,.4) .. (1.4,1.2) -- (.9,2.8);
\draw[line width=.2mm,gray] (1,2.85) .. controls (1.5,1.65) and (1.9,.6) .. (1.8,.15);
\end{tikzpicture}}}
% archive sheet number: 38
\scanpage{83}
\mspage{\handwritten{74}}
% A single diagonal cancellation runs through the main text of the sheet.
\deleted{ment constant à fibres finies et que les ${}_nG^{\unsure{o|0|\circ}}$ (pour $n\handwritten{\geq}1$) sont des schémas en groupes \deleted{\unsure{les}} finis localement libres sur $S$, que $H$ est fini localement libre sur $S$, et que \handwritten{l}a formation \deleted{\unsure{commute}} de $\underline{\operatorname{Hom}}(\underline{\unsure{\varnothing|\phi|\Phi}},\ G^{\handwritten{o\natural}})$ commute à\handwritten{\unsure{q|,}} tout changement de base. Appliquant les mêmes réflexions aux extensions de $\underline{\varnothing}_n$ par $G^o_n$ , on trouve que celles-ci correspondent aux torseurs sous le Groupe $H_n=Hx_SS_n\handwritten{\simeq}\underline{\operatorname{Hom}}(\underline{\varnothing}_n,\ G^o_n)$. Or, $H$ étant fini localement libre sur $S$, on en conclut aussitôt que la \deleted{donnée du} catégorie des torseurs \inserted{$P$} sous $H$ équivaut (par le foncteur ``complétion formelle'') à la catégorie des torseurs sous $\overset{\handwritten{\scriptscriptstyle\wedge}}{H}$ i.e. des``systêmes cohérents'' de torseurs $P_n$ sous les $H_n$ . Cela implique donc que la catégorie des extensions de $\underline{\varnothing}$ par $G^{o\handwritten{\natural}}$ équivaut à celle des \deleted{extensions de} ``systêmes cphérents'' d'extensions des $\underline{\varnothing}_n$ par les $G^o_n$ , ce qui précise donc la définition de l'extension $G^{\handwritten{\natural}}$ de (7.6.1).}

\deleted{Le schéma en groupes $G^{\handwritten{\natural}}$ , en tant qu'extension de $\underline{\varnothing}$ par $G^{o\handwritten{\natural}}$ , dé\deleted{\unsure{pe}}pend évidemment encore de façon fonctorielle de $G$. Sa formation est compatible aux changements de base par un morphisme de traits complets $S'\handwritten{\longrightarrow}S$.}

\margin{\handwritten{$\unsure{V|\checkmark}$}}
\deleted{7.6.3.Notons que les constructions \unsure{précédentes} gardent également un sens sur une base plus générale \unsure{qutun|qu'un} \deleted{trait}\inserted{sbhéma \uline{local}} complet; il suffit de partir d'un anneau noethérien adique $\deleted{\unsure{\mathbb{X}|\mathbb{A}|\mathbb{V}|W}}\handwritten{V}$, dont la topologie est définie par un \deleted{\unsure{sché}} idéal $\underline{m}\handwritten{\subset}\deleted{\mathbb{X}}\handwritten{V}$, et d'un schéma en groupes $G$ sur $\deleted{\mathbb{X}}\handwritten{V}$ tel que $G_{\unsure{o|0}}=G\otimes_{\deleted{\mathbb{X}}\handwritten{V}}\deleted{\mathbb{X}}\handwritten{V}_o$ (où $\deleted{\mathbb{X}}\handwritten{V}_o=\mathbb{X}/\underline{m}$) soit extension d'un schéma en groupes étale fini par un schéma abélien par un \deleted{schéma} tore. La question essentielle est celle de montrer que le schéma formel abélien $\overset{\handwritten{\scriptscriptstyle\wedge}}{B}$\inserted{sur $\deleted{\mathbb{X}}\handwritten{V}$} construit comme dans 7.1 en termes de $G^o$ est algébrisable; cela se voit par la méthode de 7.2 lorsqu'on suppose que $S$ est normal.}
% archive sheet number: 82
